\documentclass[../../../include/open-logic-section]{subfiles}

\begin{document}

\olfileid{his}{set}{infinitesimals}
\olsection{অতিক্ষুদ্র রাশি ও অবকলন}

সপ্তদশ শতকের শেষে নিউটন ও লাইবনিৎস
স্বতন্ত্রভাবে কলনবিদ্যা আবিষ্কার করেন।
কলনবিদ্যার একটি বিশেষ গুরুত্বপূর্ণ
প্রয়োগ হলো \emph{অবকলন}। মোটামুটি
বললে, কোনো বিন্দুতে একটি অপেক্ষকের
“পরিবর্তনের হার” বা ঢালের ধারণা
দেওয়াই অবকলনের উদ্দেশ্য।

ধারণাটি স্পষ্ট করার একটি চিত্রময়
উপায় আছে। নিচে কালো রেখায় আঁকা
$f(x) = \nicefrac{x^2}{4} + \nicefrac{1}{2}$
অপেক্ষকটি বিবেচনা করো:
\begin{center}
	\begin{tikzpicture}[scale=1]
	\draw[->, gray] (-1,0) -- (4.25,0) node[right] {$x$};
	\draw[->, gray] (0,-0.25) -- (0,5.25) node[above] {$f(x)$};

	\foreach \x/\xtext in {1/1, 2/2, 3/3, 4/4}
	\draw[shift={(\x,0)}, gray] (0pt,2pt) -- (0pt,-2pt) node[below] {$\xtext$};

	\foreach \y/\ytext in {1/1, 2/2, 3/3, 4/4, 5/5}
	\draw[shift={(0,\y)}, gray] (2pt,0pt) -- (-2pt,0pt) node[left] {$\ytext$};

	\draw[oldiagcolorC] (0.5, 9/16) -- (3.5, 57/16) -- (3.5, 9/16)--cycle;
	\draw[oldiagcolorD] (.5, 9/16) -- (2.5, 33/16) -- (2.5, 9/16) -- cycle;
	\draw[oldiagcolorE] (.5, 9/16) -- (1.5, 17/16) -- (1.5, 9/16) -- cycle;
	\draw[thick] (-1,.75) parabola bend (0,0.5) (4,4.5);
	\end{tikzpicture}
\end{center}
ধরা যাক, $c = \nicefrac{1}{2}$ বিন্দুতে
অপেক্ষকটির ঢাল বের করতে চাই। প্রথমে
একটি ত্রিভুজ আঁকি যার অতিভুজের ঢাল
ওই বিন্দুতে অপেক্ষকটির ঢালের কাছাকাছি;
যেমন, উপরের লাল ত্রিভুজটি। ত্রিভুজটির
ভূমির দৈর্ঘ্য $\beta$ হলে উচ্চতা
$f(\nicefrac{1}{2}+\beta) - f(\nicefrac{1}{2})$।
ফলে অতিভুজের ঢাল:
\[
\frac{f(\nicefrac{1}{2}+\beta) - f(\nicefrac{1}{2})}{\beta}.
\]
সুতরাং ভূমির দৈর্ঘ্য~$3$-বিশিষ্ট
আমাদের !!{colorC} ত্রিভুজটির ঢাল
ঠিক~$1$। আরও ছোট, ভূমির
দৈর্ঘ্য~$2$-বিশিষ্ট !!{colorD}
ত্রিভুজটির অতিভুজ আরও ভালো
কাছাকাছি মান দেয়; তার ঢাল
$\nicefrac{3}{4}$। আরও ছোট,
ভূমির দৈর্ঘ্য~$1$-বিশিষ্ট
!!{colorE} ত্রিভুজটির কাছাকাছি
মান আরও ভালো; তার ঢাল
$\nicefrac{1}{2}$।

ত্রিভুজগুলি ক্রমে ছোট করলে
কাছাকাছি মানও ক্রমে ভালো
হয়। তাই এমন কিছু বলা যেতে
পারে: \emph{অতিক্ষুদ্র} ভূমির
দৈর্ঘ্যবিশিষ্ট ত্রিভুজের
অতিভুজই $c = \nicefrac{1}{2}$
বিন্দুতে প্রকৃত ঢাল দেয়।
এইভাবে $c$ বিন্দুতে অপেক্ষক
$f$-এর (প্রথম) অবকলনের একটি
সূত্র পাওয়া যাবে:
\[
{f'}(c) = \frac{f(c+\beta) - f(c)}{\beta} \text{ যেখানে $\beta$ অতিক্ষুদ্র।}
\]
মোটামুটি এই কথাই নিউটন
ও লাইবনিৎস বলেছিলেন।

কিন্তু তাঁরা এ কথা বলার
পর তাঁদের জিজ্ঞাসা করতেই
হয়: \emph{অতিক্ষুদ্র রাশি}
কী? এখানে গুরুতর উভয়সংকট
দেখা দেয়। $\beta = 0$
হলে $f'$-এর সংজ্ঞা থাকে
না, কারণ তখন $0$ দিয়ে
ভাগ করতে হয়। অথচ
$\beta > 0$ হলে আমরা
প্রকৃত ঢালের বদলে শুধু
তার একটি \emph{কাছাকাছি
মান} পাই।

এই আপত্তি পরবর্তী যুগের
চাপিয়ে দেওয়া উদ্বেগ নয়।
নিউটনের অনুসারীদের সমালোচনা
করে বার্কলি বলেছিলেন:
\begin{quote}
আমি স্বীকার করি, চিহ্ন দিয়ে
কোনো বস্তু কিংবা কিছুই নয়—
দুই-ই বোঝানো যেতে পারে;
ফলে মূল সংকেতলিপি
$c + \beta$-তে $\beta$
কোনো বৃদ্ধি অথবা কিছুই
নয়, যেকোনোটি বোঝাতে
পারত। কিন্তু যেটিই
বোঝাও না কেন, তার
সেই অর্থের সঙ্গে সঙ্গতি
রেখেই যুক্তি দিতে হবে;
দুটি অর্থ একসঙ্গে ধরে
এগোনো চলবে না। তা
করলে স্পষ্ট কুযুক্তি হয়।
(\citealt[\S{}XIII]{Berkeley1734};
চলরাশিগুলি আগের আলোচনার
সঙ্গে মিলিয়ে বদলানো হয়েছে)
\end{quote}
বার্কলির আপত্তির বিরুদ্ধে
অতিক্ষুদ্র রাশির কলনবিদ্যাকে
সমর্থন করতে কেউ বলতে পারেন,
“অতিক্ষুদ্র” কথাটি কেবল
রূপক। \emph{খুব ছোট}
একটি ত্রিভুজ নিলেই
স্পর্শকের \emph{যথেষ্ট
ভালো} কাছাকাছি মান
পাওয়া যাবে। বার্কলির
এরও উত্তর ছিল: প্রকৌশলে
তা যথেষ্ট হলেও এতে
গণিতের \emph{মর্যাদা}
ক্ষুণ্ণ হয়, কারণ
\begin{quote}
আমাদের বলা হয় যে
\emph{in rebus mathematicis errores quàm minimi
non sunt contemnendi}।
[গণিতের ক্ষেত্রে ক্ষুদ্রতম
ভুলও উপেক্ষা করা চলে
না।] \citep[\S{}IX]{Berkeley1734}
\end{quote}
ইটালিক অংশটি নিউটনের
নিজের \emph{Quadrature of
Curves} (১৭০৪) থেকে
প্রায় হুবহু উদ্ধৃত।

বার্কলির দার্শনিক আপত্তি
ছিল গভীর ও তীক্ষ্ণ।
তা সত্ত্বেও কলনবিদ্যা
বিপুল সাফল্য পায় এবং
গণিতবিদেরা ভুলে না
পড়েই তা ব্যবহার করে
যান।

\end{document}
